\documentclass[../Main.tex]{subfiles}

\begin{document}

\section{Reflexion} \label{sec:Reflexion}

\subsection{Reflexion des Arbeitsprozesses}

Dieses Projekt hat etwa ein Jahr unserer Zeit in Anspruch genommen und mehrere Phasen durchlaufen. Diese greifen wir im Folgenden der Reihe nach auf und reflektieren sie.

Ganz am Anfang waren wir noch völlig unsicher und hatten ganz andere Vorstellungen davon, welches Produkt am Ende entstehen würde. Deshalb haben wir in alle Bereiche etwas hineingestochert: Wir haben mit Spielsystemen angefangen, von denen wir ohne Recherche annahmen, dass wir sie brauchen würden, haben uns überlegt, wie wir unsere Erkenntnisfrage beantworten wollen, und haben gelernt, wie man recherchiert. Bis zur Abgabe hatten wir noch fast ein Jahr Zeit, bis zur Zwischenpräsentation am 17.~Juni~2026 noch \textcolor{red}{[x]}~Monate. Wir waren noch wenig zielorientiert und machten vor allem das, was uns gefiel. Das mag planlos wirken, rückblickend war es das aber nicht: Wir haben ausgekundschaftet, was es gibt. Zudem stand in dieser Phase noch wenig auf dem Spiel, denn egal, was wir machten, nichts hätte ernsthaft verworfen werden müssen, und die einzelnen Arbeitsschritte hingen noch kaum voneinander ab. Dass das gut ging, war allerdings eher ein glücklicher Zufall als strukturiertes Denken unsererseits, doch so verläuft der Anfang eines Projekts oft. Könnten wir noch einmal beginnen, würden wir es genauso angehen, mit dem einzigen Unterschied, dass wir uns bewusst wären, in dieser Phase noch eine Narrenfrist zu haben.

\begin{itemize}[leftmargin=1.5cm]
    \item \textcolor{red}{[x] einsetzen: ab dem PU-Start im November 2025 etwa sieben Monate, ab Beginn der Hauptentwicklung im Februar 2026 (laut Git) etwa vier.}
\end{itemize}

Inzwischen hatten wir mit unserer Betreuungsperson alles Nötige vertraglich festgehalten. Bis zur Zwischenpräsentation spielte die Beantwortung der Erkenntnisfrage jedoch eine untergeordnete Rolle, denn die Systeme, die für die Playtests und damit für die Erkenntnisfrage relevant sind, brauchten eine technische Grundlage, die wir noch nicht hatten. Wir arbeiteten also weiter an diesen Voraussetzungen: am Spielercharakter mit seinen Animationen und seiner Steuerung, an den Gegnern und an allem, was es braucht, damit das Spiel überhaupt startet. Rückblickend war das eine gute Entscheidung, auch wenn wir die Recherche dafür etwas vernachlässigt haben. Wir haben so entschieden, weil uns dieser Teil am meisten Freude bereitete, aber auch, weil wir zuerst eine Vorstellung von unserem eigenen Spiel entwickeln mussten. Dabei war die gegenseitige Kommunikation besonders wichtig. Erst nachdem wir die ersten Systeme gebaut hatten, wussten wir, worauf wir in anderen Spielen achten mussten, und konnten mit diesem Wissen eine gute Recherche machen. Bei einem nächsten Projekt würden wir die Reihenfolge davon abhängig machen, wie konkret unsere Vorstellungen und wie gross unser Vorwissen sind: Haben wir beides, recherchieren wir zuerst; fehlt eines davon, sammeln wir zuerst eigene Erfahrungen.

Gleichzeitig rückte die Zwischenpräsentation näher, und wir mussten zusehen, dass die grundlegenden Systeme fertig wurden. Schon jetzt mussten wir Features, die wir zwar wollten, die für die Erkenntnisfrage aber nicht besonders relevant waren, streichen oder stark eingrenzen und priorisieren, was für ein Spiel essenziell ist. Dazu kam, dass wir bis dahin einen begründet geplanten Playtest durchgeführt haben mussten. Das Zeitmanagement spielte deshalb eine grosse Rolle, denn die Checkliste aus dem Vertrag hatte noch viele offene Punkte. Bis zur Zwischenpräsentation haben wir alles erreicht, was wir vertraglich festgehalten hatten, auch wenn es sehr knapp wurde. Nicht das Fertigstellen, sondern vor allem die Qualitätssicherung hat uns dabei einen Strich durch die Rechnung gemacht, weil wir möglichst wenige Kompromisse eingehen und eine möglichst hohe Qualität erreichen wollten. Im Nachhinein wäre es interessant gewesen zu erfahren, wie eine Arbeit mit nur 80 oder 85 Prozent unseres Anspruchs bewertet worden wäre. Da die Zwischenbewertung nicht in die Note der Maturitätsarbeit einfliesst, haben wir hier eine Chance verpasst. Andererseits konnten wir Erfahrung mit Crunchen sammeln, die es uns am Ende ermöglicht hat, mit genügend Zeit und kaum Crunchen 95 bis 100 Prozent herauszuholen.

\begin{itemize}[leftmargin=1.5cm]
    \item \textcolor{red}{\enquote{Crunch} ist Fachjargon aus der Spieleentwicklung: im Glossar erklären oder umschreiben.}
\end{itemize}

Die Rückmeldung zur Zwischenpräsentation war in vielem positiv, in einigem weniger. Die Kritik betraf weniger den Prototyp als die begleitenden Texte, die einige Schwachstellen hatten. Beim Umsetzen dieser Kritik hatten wir zunächst Schwierigkeiten, weil wir die bemängelten Punkte nicht verinnerlicht hatten -- sonst hätten wir sie von Anfang an richtig gemacht. Deshalb stellten wir unserer Betreuungsperson so lange Fragen, bis wir sicher waren, sie verstanden zu haben.

Nach der Zwischenpräsentation hatten wir wieder genug Zeit, und die Abgabe war weit genug entfernt, dass wir unsere Ambitionen beim Projektumfang vergrösserten. Gleichzeitig nahm das Tempo, mit dem wir Neues produzierten, ab. Das lag zum einen daran, dass das Projekt inzwischen so gross war, dass kleinere Neuerungen wenig Unterschied machten. Zum anderen war unser anfänglicher Eifer zwar noch lange nicht erloschen, aber das Projekt zu pflegen und auch Aufgaben zu erledigen, die uns weniger Spass machten und sich mittlerweile angehäuft hatten, setzte uns zunehmend zu.

\begin{itemize}[leftmargin=1.5cm]
    \item \textcolor{red}{\textbf{1. Ausgangslage und Plan:} Erste Commits im Juni 2025; der Vertrag sah vor: Zwischenprodukt 10.6., Zwischenpräsentation 17.6., Analyse der Referenzspiele Anfang Juli, Spielsysteme Mitte August, Doku und Fragebögen Ende September, Playtests Mitte Oktober, Verbesserungen Ende Oktober.}
    \item \textcolor{red}{\textbf{2. Verlauf in Phasen} (aus Git und Journal): Hauptentwicklung des Prototyps Februar bis Juni 2026; Juni: README und Architekturdiagramm fürs Zwischenprodukt; 16.--20.7.: Kapitel 1--4 und die \game{NieR}-Analyse als Entwurf; Juli bis September kaum Arbeit am Spiel (Sommerpause?); Oktober: Level überarbeitet, Playtests vorbereitet und durchgeführt; November: Begleitbericht fertiggestellt.}
    \item \textcolor{red}{\textbf{3. Scope-Entscheide:} 20.5.: Vertrag bereinigt, Projektumfang reduziert, Erkenntnisfrage festgelegt -- was wurde gestrichen? Später: Bonus-Inhalte weggelassen, technische Dokumentation ausgelagert, README gestrichen, Recherche in Kapitel~\ref{sec:Leveldesign} auf Rat der Betreuungsperson auf das für die Erkenntnisfrage Nötige gekürzt.}
    \item \textcolor{red}{\textbf{4. Plan und tatsächlicher Verlauf:} Playtests 26.10.--2.11.\ statt Mitte Oktober, Verbesserungen erst im Bericht statt Ende Oktober. Warum?}
    \item \textcolor{red}{\textbf{5. Wichtige Wendepunkte} (Journal 24./25.10., 13~h): Obwohl der Prototyp technisch fertig war, machten die Kämpfe bei eigenen Tests wenig Spass. Ihr habt die Encounters gekürzt, die Gegner aggressiver gemacht und Rechercheergebnisse nachgeholt, die noch nicht umgesetzt waren. Eure eigene Erkenntnis dort: \enquote{Die technische Fertigstellung allein garantiert noch kein gutes Spielerlebnis} -- und weil ihr das Level so oft gespielt habt, braucht es externe Playtests. Das ist ein starkes Reflexionsbeispiel.}
    \item \textcolor{red}{\textbf{6. Rückmeldungen der Zwischenbewertung und was ihr daraus gemacht habt:} Zeitaufwand konsequenter dokumentieren, Reflexionen vertiefen, Entscheidungen mit Alternativen begründen, KI vollständig dokumentieren, Planung mit Arbeitspaketen, Bericht und Code verknüpfen.}
    \item \textcolor{red}{\textbf{7. Was gut funktioniert hat:} z.\,B. iteratives Vorgehen, Git und GitHub, Playtest-Vorgehen -- jeweils mit Grund.}
    \item \textcolor{red}{\textbf{8. Methodik im Rückblick:} kurz, Einzelheiten in den Limitationen (Abschnitt~\ref{sec:Limitationen}).}
    \item \textcolor{red}{\textbf{9. Was ihr anders machen würdet}, z.\,B. früher extern testen, Rechercheergebnisse von Anfang an mit dem Level abgleichen, Bericht laufend statt am Schluss schreiben.}
    \item \textcolor{red}{\textbf{10. Verweis aufs Arbeitsjournal}, wo der Verlauf im Detail dokumentiert ist.}
\end{itemize}

\subsection{Gruppenarbeit und Aufgabenteilung}

\begin{itemize}[leftmargin=1.5cm]
    \item \textcolor{red}{\textbf{1. Wie ihr die Arbeit aufgeteilt habt.} Laut Zwischenbewertung: Rodin stärker bei Animation, UI, Bericht und README; Zhile stärker bei Architektur, Bewegung, State Machines und Kampf-Grundlagen. Das deckt sich mit Git: Zhile vor allem \code{Entities/Player}, \code{Common} (Utility, Tools), \code{Core}, Kamera; Rodin vor allem \code{Interfaces} und Menüs sowie die TrenchBroom-Karten des Levels.}
    \item \textcolor{red}{\textbf{2. Tabelle Aufgabenteilung und Eigenleistung} für Prototyp und Bericht (Rodin / Zhile / gemeinsam). Zum Bericht: Kapitel 1--8 und die \game{NieR}-Analyse Rodin, die \game{Ninja-Gaiden}-Analyse Zhile, README und Architekturdiagramm (Juni) vor allem Zhile. Bitte prüfen und ergänzen (Playtests, Leveldesign, Audio, \ldots).}
    \item \textcolor{red}{\textbf{3. Zusammenarbeit und Kommunikation:} Wie habt ihr euch abgesprochen (Treffen, Discord, GitHub)? Wie habt ihr Entscheidungen getroffen?}
    \item \textcolor{red}{\textbf{4. Veränderung im Verlauf:} z.\,B. ab Oktober Rodin vor allem Level, Playtests und Bericht, Zhile weiter am Spiel (laut Git im Oktober sehr aktiv). Wie kam es dazu?}
    \item \textcolor{red}{\textbf{5. Fazit:} Was lief gut, was war schwierig, was würdet ihr in einer Gruppenarbeit anders machen? Hinweis: Commit-Zahlen sagen wenig über den Aufwand aus, deshalb nicht als Massstab verwenden.}
\end{itemize}

\subsection{Verwendung von KI} \label{sec:ReflexionKI}

\begin{itemize}[leftmargin=1.5cm]
    \item \textcolor{red}{\textbf{1. Welche Tools wofür:} Code (Treppen-Werkzeug und Animations-Hilfsskript im Frühling, Audiosystem im Oktober), Begleitbericht (Korrektur, Struktur, Rückmeldungen). Details in der KI-Deklaration (Abschnitt~\ref{sec:KIDeklaration}).}
    \item \textcolor{red}{\textbf{2. Was geholfen hat:} z.\,B. Widersprüche zwischen Kapiteln finden, unbelegte Aussagen markieren, \LaTeX-Formatierung, Fehlersuche im Code.}
    \item \textcolor{red}{\textbf{3. Grenzen und Probleme:} Die KI ging mehrmals von falschen Annahmen aus, die ihr korrigieren musstet, z.\,B. dass Themenfindung mit KI gemacht wurde, dass ein Tutorial geplant war, welche Räume Brückenräume sind, oder dass in den Protokollen ein anderer Archer gemeint war.}
    \item \textcolor{red}{\textbf{4. Wie ihr die Ergebnisse überprüft habt:} Abgleich mit Quellen, Playtest-Protokollen und dem Spiel selbst; Vorschläge nicht ungeprüft übernommen, einige bewusst verworfen.}
    \item \textcolor{red}{\textbf{5. Eigenanteil sichern:} Interpretation und Urteile selbst festgelegt; von der KI formulierte Stellen markiert und in eigenen Worten neu geschrieben; Regel \enquote{nicht wortwörtlich übernehmen} aus dem Merkblatt.}
    \item \textcolor{red}{\textbf{6. Fazit:} Würdet ihr KI wieder so einsetzen? Was würdet ihr anders machen, z.\,B. KI-Nutzung von Anfang an laufend im Journal dokumentieren?}
\end{itemize}

\subsection{Persönliche Erkenntnisse}

\subsubsection{Rodin Arslan}

\begin{itemize}[leftmargin=1.5cm]
    \item \textcolor{red}{\textbf{1. Fachlich:} z.\,B. Animation (u.\,a. Hängeanimation), UI und Menüs, Leveldesign mit TrenchBroom, Gegner- und Encounter-Design, Playtesting, \LaTeX\ und BibTeX.}
    \item \textcolor{red}{\textbf{2. Methodisch:} Spielanalyse mit dem MDA-Framework, Recherche und Zitieren, Playtests planen und auswerten, Projektplanung.}
    \item \textcolor{red}{\textbf{3. Persönlich:} z.\,B. dass ein technisch fertiges Spiel noch nicht gut spielbar ist (Journal 24./25.10.), Zeitmanagement in der Schlussphase, Umgang mit Rückmeldungen.}
    \item \textcolor{red}{\textbf{4. Was du mitnimmst} für die Zukunft.}
\end{itemize}

\subsubsection{Zhile Zhu}

\textcolor{red}{ZHILE DAS IST FÜR DICH}

\end{document}
